% Resumen ejecutivo de dos páginas para el comité de dirección (apartado 9.4 del
% enunciado). Mismo estilo que docs/report. Las cifras proceden de la ejecución de
% referencia del 10 de octubre de 2026 (results/escenarios/tablas/).
\documentclass[10pt,a4paper]{article}

\usepackage[T1]{fontenc}
\usepackage[utf8]{inputenc}
\usepackage[spanish,es-tabla,es-nodecimaldot,es-noshorthands]{babel}
\spanishplainpercent
\usepackage{lmodern}
\usepackage{microtype}
\usepackage[a4paper,left=17mm,right=17mm,top=14mm,bottom=14mm,headheight=12pt,headsep=5mm,footskip=7mm]{geometry}
\usepackage{graphicx}
\graphicspath{{../report/figuras/}}
\usepackage[table]{xcolor}
\definecolor{azul}{HTML}{0072B2}
\definecolor{optimista}{HTML}{009E73}
\definecolor{pesimista}{HTML}{D55E00}
\definecolor{tinta}{HTML}{252A30}
\definecolor{gris}{HTML}{60666D}
\definecolor{fondo}{HTML}{F1F5F8}
\color{tinta}
\usepackage{booktabs,tabularx,array,makecell}
\renewcommand{\arraystretch}{1.1}
\newcolumntype{Y}{>{\raggedright\arraybackslash}X}
\newcolumntype{R}{>{\raggedleft\arraybackslash}X}
\newcolumntype{N}{>{\raggedleft\arraybackslash}p{28mm}}
\usepackage{eurosym}
\usepackage{enumitem}
\setlist{nosep,leftmargin=1.2em,topsep=2pt}
\usepackage{caption}
\captionsetup{font=small,labelfont=bf,labelsep=period,justification=justified,singlelinecheck=false,skip=3pt}
\usepackage{titlesec}
\titleformat{\section}{\large\bfseries}{\thesection.}{.5em}{}
\titlespacing*{\section}{0pt}{7pt plus 2pt}{3pt}
\usepackage{fancyhdr}
\pagestyle{fancy}
\fancyhf{}
\fancyhead[L]{\small\color{gris}Kaido Motors España · Resumen ejecutivo A013}
\fancyhead[R]{\small\thepage/2}
\renewcommand{\headrulewidth}{0pt}
\usepackage[skins]{tcolorbox}
\newtcolorbox{idea}[1]{enhanced,colback=fondo,colframe=gris!35,
  boxrule=.4pt,arc=0pt,left=7pt,right=7pt,top=4pt,bottom=4pt,
  title={#1},colbacktitle=fondo,coltitle=tinta,fonttitle=\bfseries,
  titlerule=0pt,before skip=5pt,after skip=5pt}
\usepackage[unicode,colorlinks=true,linkcolor=azul,urlcolor=azul,
  pdftitle={Kaido Motors España: resumen ejecutivo del plan A013},
  pdfauthor={Proyecto Kaido Motors},pdfsubject={Simulación de Sistemas Logísticos}]{hyperref}
\newcommand{\MEE}{\,M\euro}
\newcommand{\fuente}[1]{\par\smallskip{\footnotesize\color{gris}\textit{Fuente:} #1.\par}}
\newcommand{\pes}{\textcolor{pesimista}{\textbf{Pesimista}}}
\newcommand{\cen}{\textcolor{azul}{\textbf{Central}}}
\newcommand{\opt}{\textcolor{optimista}{\textbf{Optimista}}}
\setlength{\parindent}{0pt}
\setlength{\parskip}{4pt}
\setlength{\tabcolsep}{4.2pt}

\begin{document}
\thispagestyle{empty}
\noindent
\begin{minipage}[b]{.74\linewidth}
  {\scshape\small Simulación de Sistemas Logísticos · Para el comité de dirección}\par
  \vspace{1mm}
  {\Large\bfseries Kaido Motors España: plan de fabricación A013}\par
  \vspace{.5mm}
  {\large\color{azul}Resumen ejecutivo}
\end{minipage}\hfill
\begin{minipage}[b]{.2\linewidth}\raggedleft
  \includegraphics[width=\linewidth]{logo-uax.png}
\end{minipage}
\par\vspace{1mm}{\color{gris!40}\rule{\linewidth}{.4pt}}

\begin{idea}{La propuesta}
Se propone fabricar \textbf{16.144 coches} en A013, con el reparto mensual de la
tabla~\ref{tab:plan}, y corregir los pedidos cada mes con las ventas reales. En el
escenario central, este plan vende unos \textbf{18.540 coches}, atiende al
\textbf{97,6\,\%} de los compradores, deja unos 2.220 coches en la red a fin de año
y obtiene \textbf{62,9\MEE} después de pagar la venta y el stock. MSUV, SSUV y CSUV
aportan el 93\,\% de ese resultado, por lo que conviene darles prioridad. C no
necesita pedidos hasta agosto porque empieza el año con mucho stock, y B deja muy
poco margen: si la matriz carga la prima de CO$_2$, cada coche adicional de B
genera pérdidas.
\end{idea}

\section{Cuántos coches fabricar cada mes}
Lo que se pide en un mes llega a la red tres meses después: lo que se venda
hasta marzo ya viene en camino. La tabla recoge los pedidos que se deciden en
A013; los de octubre a diciembre son coches para enero--marzo de A014.

\begin{table}[h]\centering
\caption{Plan inicial de fabricación por mes de pedido, en coches. Escenario central.}\label{tab:plan}
\small
\begin{tabular}{@{}lrrrrrrrrrrrr>{\bfseries}r@{}}
\toprule
Segmento & Ene & Feb & Mar & Abr & May & Jun & Jul & Ago & Sep & Oct & Nov & Dic & Total \\
\midrule
MSUV · SUV medio    & 353 & 622 & 693 & 593 & 293 & 336 & 452 & 516 & 633 & 314 & 309 & 468 & 5.582 \\
SSUV · SUV pequeño  & 659 & 494 & 750 & 530 & 253 & 245 & 325 & 412 & 496 & 219 & 335 & 426 & 5.144 \\
CSUV · SUV compacto & 0 & 218 & 347 & 310 & 149 & 112 & 191 & 253 & 272 & 122 & 185 & 257 & 2.416 \\
C · compacto        & 0 & 0 & 0 & 0 & 0 & 0 & 0 & 172 & 220 & 163 & 201 & 232 & 988 \\
B · pequeño         & 153 & 132 & 168 & 121 & 66 & 45 & 67 & 78 & 100 & 64 & 103 & 136 & 1.233 \\
MD · berlina media  & 0 & 15 & 59 & 36 & 17 & 29 & 42 & 33 & 33 & 31 & 40 & 45 & 380 \\
S · descapotable    & 80 & 72 & 74 & 30 & 0 & 0 & 0 & 0 & 0 & 56 & 83 & 6 & 401 \\
\midrule
\rowcolor{fondo}
\textbf{Total} & 1.245 & 1.553 & 2.091 & 1.620 & 778 & 767 & 1.077 & 1.464 & 1.754 & 969 & 1.256 & 1.570 & 16.144 \\
\bottomrule
\end{tabular}
\fuente{\texttt{results/escenarios/tablas/plan\_pedidos\_inicial\_a013\_a014.csv}}
\end{table}

\vspace{-2mm}
Cada pedido cubre la demanda prevista de su mes de llegada, más un colchón por
segmento para atender al menos al 95\,\% de los compradores en el escenario
central, menos el stock que habrá en la red. En cada segmento se eligió el
colchón que más resultado deja cumpliendo ese 95\,\%. El plan es el punto de partida: al revisar cada
mes, el central acaba fabricando de media 15.756 coches.

Tres datos explican la forma del plan:
\begin{itemize}
\item \textbf{C} empieza el año con mucho stock: se venden los coches que ya hay
  y no se piden más hasta agosto.
\item \textbf{S} no recibe pedidos de mayo a septiembre: vende unos 370 coches al
  año y su stock ya cubre meses de venta.
\item \textbf{Mayo y junio} bajan porque llegan en agosto y septiembre, meses de
  poca venta; marzo sube para llegar en junio, el mes de más demanda.
\end{itemize}

\section{Qué pasa con el plan en el escenario central}
Se venden 18.539 coches de una demanda de 19.017; se pierden unos 480
compradores, el 2,4\,\%. El stock medio en la red es de 4.075 coches, unos
2,6 meses de venta, y a fin de año quedan 547 coches con más de 90 días.
El gasto comercial (VM) es de 40,5\MEE{} y el stock le cuesta a la marca
2,7\MEE. Queda un margen (CVP) de 65,6\MEE{} y un resultado de 62,9\MEE{}.
El plan no recurre a transfers, los coches traídos del stock de otros
países europeos. En una de cada diez repeticiones el resultado baja de 56,4\MEE{} y en
otra de cada diez supera los 69,5\MEE.

\clearpage
\section{Cuánto cambia si el año sale peor o mejor}
Las tres columnas parten del mismo plan. En pesimista y optimista la demanda y los
costes son otros, y los pedidos solo se pueden corregir a tres meses vista.

\begin{table}[h]\centering
\caption{Resultados de A013 con el plan propuesto (tabla 5 del enunciado), total
de la marca. Medias de 1.000 repeticiones por escenario; entre paréntesis, la
diferencia con el central.}\label{tab:tabla5}
\small
\begin{tabularx}{\linewidth}{@{}YNNN@{}}
\toprule
Indicador & \pes & \cen & \opt \\
\midrule
Demanda de clientes (coches) & 15.652 & 19.017 & 28.445 \\
Coches vendidos (coches) & 12.151 \,{\color{gris}\footnotesize(\textminus6.388)} & 18.539 & 24.952 \,{\color{gris}\footnotesize(+6.413)} \\
Coches fabricados, pedidos A013 (coches) & 6.371 \,{\color{gris}\footnotesize(\textminus9.385)} & 15.756 & 23.268 \,{\color{gris}\footnotesize(+7.512)} \\
Stock medio en la red (coches; meses de venta) & 3.489; 2,7 & 4.075; 2,6 & 3.358; 1,4 \\
Coches con más de 90 días a fin de año & 601 & 547 & 404 \\
Coste del stock para la marca (M\euro) & 2,20 & 2,73 & 2,38 \\
Ventas perdidas (coches) & 3.501 & 478 & 3.493 \\
\quad Margen no obtenido por ellas (M\euro) & \textminus3,09 & 1,47 & 12,96 \\
Transfers (coches y M\euro) & 0 & 0 & 0 \\
Gasto comercial, VM (M\euro) & 27,63 \,{\color{gris}\footnotesize(\textminus12,87)} & 40,50 & 53,36 \,{\color{gris}\footnotesize(+12,86)} \\
Margen de la marca, CVP (M\euro) & 46,49 \,{\color{gris}\footnotesize(\textminus19,12)} & 65,62 & 88,23 \,{\color{gris}\footnotesize(+22,61)} \\
\midrule
\quad Prima de CO$_2$ / arancel / tipo de cambio (M\euro) & 36,56 / 4,64 / 0 & 0 / 0 / 0 & 0 / 0 / \textminus5,52 \\
\rowcolor{fondo}
\textbf{Resultado tras costes (M\euro)} & \textbf{3,09} \,{\color{gris}\footnotesize(\textminus59,79)} & \textbf{62,88} & \textbf{91,37} \,{\color{gris}\footnotesize(+28,49)} \\
\quad Intervalo P10--P90 (M\euro) & 1,87--4,35 & 56,45--69,47 & 84,52--98,39 \\
Compradores atendidos (\%) & 77,8 & 97,6 & 87,9 \\
\bottomrule
\end{tabularx}
\fuente{\texttt{results/escenarios/tablas/tabla5\_plan\_recomendado\_a013.csv}.
Resultado = CVP $-$ stock $-$ transfers $-$ CO$_2$ $-$ arancel $-$ cambio; un coste
de cambio negativo es un ahorro. Detalle por segmento en el anexo del informe}
\end{table}

\vspace{-2mm}
\textbf{Pesimista.} El resultado baja a 3,1\MEE, en parte por la menor demanda
y sobre todo por dos costes nuevos: la prima de CO$_2$ (36,6\MEE) y un arancel
(4,6\MEE). Con ellos, cada coche nuevo de B, C, SSUV y S pierde dinero, de modo que
el plan deja de pedirlos y vende el stock que ya hay en la red. Así quedan sin
atender unos 3.500 compradores, cuyas ventas habrían restado 3,1\MEE. Si se
siguiera pidiendo para atenderlos, el servicio subiría al 98,6\,\% y el resultado
sería de \textminus0,6\MEE; con un arancel de 1.000\,\euro{} por coche, el año
terminaría en \textminus1,5\MEE.

\textbf{Optimista.} Se venden unos 6.400 coches más y el resultado sube a
91,4\MEE, de los que 5,5\MEE{} proceden del ahorro por tipo de cambio. La demanda
crece más deprisa de lo que permiten corregir los pedidos, así que quedan sin
atender unos 3.490 compradores, sobre todo de SSUV y MSUV, con 13,0\MEE{} de margen
no obtenido. El stock medio baja a 1,4 meses de venta.

\begin{table}[h]\centering
\caption{Lo esencial por segmento: medias de A013 con el plan propuesto.}\label{tab:segmentos}
\small
\begin{tabularx}{\linewidth}{@{}l*{9}{R}@{}}
\toprule
& \multicolumn{3}{c}{Coches vendidos} & \multicolumn{3}{c}{Coches fabricados (A013)} & \multicolumn{3}{c}{Resultado (M\euro)} \\
\cmidrule(lr){2-4}\cmidrule(lr){5-7}\cmidrule(l){8-10}
Segmento & \textcolor{pesimista}{Pes.} & \textcolor{azul}{Cen.} & \textcolor{optimista}{Opt.}
         & \textcolor{pesimista}{Pes.} & \textcolor{azul}{Cen.} & \textcolor{optimista}{Opt.}
         & \textcolor{pesimista}{Pes.} & \textcolor{azul}{Cen.} & \textcolor{optimista}{Opt.} \\
\midrule
MSUV & 4.556 & 5.554 & 7.067 & 4.335 & 5.501 & 7.292 & 4,72 & 31,91 & 42,61 \\
SSUV & 1.915 & 5.387 & 7.077 & 0 & 4.882 & 6.822 & \textminus1,03 & 14,86 & 21,60 \\
CSUV & 2.816 & 3.327 & 4.692 & 1.751 & 2.357 & 3.935 & 3,52 & 11,77 & 18,01 \\
C    & 1.796 & 2.258 & 3.451 & 0 & 959 & 2.283 & \textminus3,16 & 1,15 & 4,00 \\
B    & 426 & 1.212 & 1.578 & 0 & 1.170 & 1.602 & \textminus1,24 & 0,02 & 0,47 \\
MD   & 360 & 433 & 571 & 286 & 375 & 546 & 0,45 & 1,63 & 2,31 \\
S    & 282 & 368 & 516 & 0 & 513 & 788 & \textminus0,16 & 1,55 & 2,37 \\
\midrule
\rowcolor{fondo}
\textbf{Total} & 12.151 & 18.539 & 24.952 & 6.371 & 15.756 & 23.268 & 3,09 & 62,88 & 91,37 \\
\bottomrule
\end{tabularx}
\fuente{\texttt{results/escenarios/tablas/tabla5\_plan\_recomendado\_a013.csv}}
\end{table}

\vspace{-2mm}
\section{Qué decidir y qué vigilar}
\begin{enumerate}
\item \textbf{Aprobar el plan de la tabla~\ref{tab:plan} como punto de partida} y
  revisarlo cada mes: solo se puede cambiar lo que llega dentro de tres meses o más.
\item \textbf{Dar prioridad a MSUV, SSUV y CSUV.} Son los que más dejan y los que
  más compradores pierden si el año sale mejor. SSUV es el más ajustado: en cuatro
  de cada diez repeticiones del central atiende a menos del 95\,\%.
\item \textbf{No pedir coches que pierden dinero.} Si se confirma la prima de
  CO$_2$ o el arancel, frenar los pedidos nuevos de B, C, SSUV y S aunque haya
  demanda, y vender primero lo que ya está en la red.
\item \textbf{Vigilar las ventas de tres meses frente a la previsión}, segmento a
  segmento. Si las ventas caen cuando la red está vacía, la causa es la falta de
  coches y no una menor demanda.
  CSUV, con poca historia, es el segmento con la previsión más incierta.
\end{enumerate}

{\footnotesize\color{gris}Los escenarios no tienen probabilidad asignada; el
intervalo P10--P90 recoge la variación dentro de cada uno. El resultado es la
contribución tras los costes del modelo, no el beneficio contable. Datos y nombres
camuflados; material de uso exclusivamente docente.\par}
\end{document}
